\subsubsection{Refinamiento del módulo y conservación de los dos lectores}
\label{mg03:refinamiento}

La identidad \eqref{eq:defecto-integrado} conserva el defecto entero
\(1215\) mediante el residuo \(54\) y el cociente \(129\). Su
generalización requiere declarar por separado el módulo de lectura,
el tamaño del soporte y el número de coordenadas. Estas tres operaciones
producen leyes de conteo distintas a partir de las mismas hojas APP.

Para \(N\geq2\), fijamos el representante positivo
\[
 \rho_N(s)=1+((s-1)\bmod N)\in\{1,\ldots,N\},
 \qquad s=\rho_N(s)+NQ_N(s),\quad Q_N(s)\in\mathbb Z.
\]
Sobre \(\{1,\ldots,N\}^2\), las lecturas plenas son
\(\rho_N(i+j)\) y \(\rho_N(ij)\). Sus sumas se denotan
\(S_+^{\mathrm{full}}(N)\) y \(S_\times^{\mathrm{full}}(N)\).

\begin{proposition}[Ley del defecto para un módulo entero]
Para todo \(N\geq2\),
\begin{align}
 S_+^{\mathrm{full}}(N)&=\frac{N^2(N+1)}2,
 \label{mg03:suma-plena}\\
 S_\times^{\mathrm{full}}(N)&=
 \frac N2\left(N^2+\sum_{i=1}^{N}\gcd(i,N)\right),
 \label{mg03:producto-pleno}\\
 \Delta_N^{\mathrm{full}}&=
 \frac N2\left(\sum_{i=1}^{N}\gcd(i,N)-N\right).
 \label{mg03:defecto-pleno}
\end{align}
\end{proposition}

\begin{proof}
En una fila aditiva, la traslación por \(i\) permuta las \(N\)
clases y sus representantes positivos. La suma por fila es
\(N(N+1)/2\), lo que demuestra \eqref{mg03:suma-plena}.
Para una fila multiplicativa, sea \(g=\gcd(i,N)\). La imagen de
\(j\mapsto ij\pmod N\) contiene las \(N/g\) clases múltiplos de
\(g\), y cada una tiene \(g\) preimágenes. Los representantes
positivos son \(g,2g,\ldots,N\), cuya suma es
\[
 g\frac{(N/g)(N/g+1)}2=\frac{N(N+g)}{2g}.
\]
La multiplicidad \(g\) da \(N(N+g)/2\) por fila. Al sumar sobre
\(i\) se obtiene \eqref{mg03:producto-pleno}; su diferencia con
\eqref{mg03:suma-plena} produce \eqref{mg03:defecto-pleno}.
\end{proof}

\begin{corollary}[Módulos primo-potencia]
Si \(N=p^m\), con \(p\) primo y \(m\geq1\), entonces
\begin{equation}
 \sum_{i=1}^{p^m}\gcd(i,p^m)
 =p^m+m(p-1)p^{m-1},\qquad
 \Delta_{p^m}^{\mathrm{full}}
 =\frac{m(p-1)}2p^{2m-1}.
 \label{mg03:primo-potencia}
\end{equation}
\end{corollary}

\begin{proof}
Para cada \(0\leq k<m\), hay
\(p^{m-k}-p^{m-k-1}\) índices de valoración \(p\)-ádica
exactamente \(k\). Su máximo divisor común con \(p^m\) vale
\(p^k\), y la contribución de ese estrato es
\((p-1)p^{m-1}\). Los \(m\) estratos aportan
\(m(p-1)p^{m-1}\), mientras que el índice \(i=p^m\) aporta
\(p^m\). La sustitución en \eqref{mg03:defecto-pleno} da la
segunda fórmula.
\end{proof}

La especialización triádica es
\(\Delta_{3^m}^{\mathrm{full}}=m3^{2m-1}\). Los primeros módulos
muestran simultáneamente las dos sumas y su diferencia:
\[
\begin{array}{c|r|r|r}
 N&S_+^{\mathrm{full}}&S_\times^{\mathrm{full}}&\Delta_N^{\mathrm{full}}\\\hline
 3&18&21&3\\
 9&405&459&54\\
 27&10206&10935&729\\
 81&269001&277749&8748
\end{array}
\]

El refinamiento periódico mantiene, en cambio, \(\rho_9\) y amplía
el soporte a \(\{1,\ldots,9r\}^2\), con \(r\geq1\) entero. Cada par de residuos módulo
nueve tiene \(r^2\) representantes en ese soporte, tanto para la
suma como para el producto. Por consiguiente,
\begin{equation}
 S_+^{\mathrm{per}}(9r;9)=r^2\,405,\quad
 S_\times^{\mathrm{per}}(9r;9)=r^2\,459,\quad
 \Delta^{\mathrm{per}}(9r;9)=r^2\,54.
 \label{mg03:periodico}
\end{equation}
Para \(r=3\), las sumas valen \(3645\) y \(4131\), y el defecto
vale \(486\). La lectura plena en \(27\) produce \(729\) porque
actualiza tanto el módulo como sus representantes; la lectura periódica
produce \(486\) porque conserva el módulo nueve y multiplica sus
incidencias por nueve. En ambas lecturas, las divisiones
\(i+j=R^++NQ^+\), \(ij=R^\times+NQ^\times\), con el módulo
efectivamente empleado, mantienen la identidad entera por celda y su
suma. Así queda fijado el dato que debe acompañar a cada defecto.

\subsubsection{Variación dimensional con módulo nueve fijo}
\label{mg03:dimension}

Una segunda extensión mantiene el módulo y aumenta el número de
factores. Para \(d\geq1\), sean
\[
 \mathcal A_d=\{1,\ldots,9\}^d,\qquad
 \Sigma_d(x)=\rho_9\!\left(\sum_{j=1}^d x_j\right),\qquad
 \Pi_d(x)=\rho_9\!\left(\prod_{j=1}^d x_j\right).
\]

\begin{proposition}[Conteo dimensional de acumulación y producto]
Las sumas de ambos lectores sobre \(\mathcal A_d\) son
\begin{equation}
 S_{+,d}=45\,9^{d-1},\qquad
 S_{\times,d}=9^{d+1}-9(d+3)6^{d-1}.
 \label{mg03:ley-dimensional}
\end{equation}
Por tanto,
\begin{equation}
 \frac{S_{\times,d}-S_{+,d}}{9^d}
 =4-(d+3)\left(\frac23\right)^{d-1}
 \longrightarrow4.
 \label{mg03:limite-dimensional}
\end{equation}
\end{proposition}

\begin{proof}
Fijadas las primeras \(d-1\) coordenadas, la última recorre una vez
todos los residuos de la suma. Cada representante positivo aparece
\(9^{d-1}\) veces, y su suma es \(45\); esto prueba la fórmula
aditiva.

Para el producto, se distinguen las unidades
\(U=\{1,2,4,5,7,8\}\), los múltiplos de tres de valoración uno
\(T=\{3,6\}\), y el representante nulo \(9\). Las \(6^d\)
palabras con todos sus factores en \(U\) distribuyen sus productos
uniformemente entre las seis unidades: al fijar \(d-1\) factores,
el último actúa por una permutación del grupo. Su contribución es
\(27\,6^{d-1}\).

Las palabras con exactamente un factor en \(T\) y los demás en
\(U\) son \(2d6^{d-1}\). Para cada posición del factor singular y
cada elección de las unidades restantes, las dos elecciones \(3,6\)
producen una vez cada uno de los residuos \(3,6\). La contribución
es \(9d6^{d-1}\). Las palabras restantes contienen un factor
\(9\), o al menos dos factores en \(T\), y su producto tiene
representante \(9\). Su número es
\(9^d-6^d-2d6^{d-1}\). En consecuencia,
\[
 S_{\times,d}=27\,6^{d-1}+9d6^{d-1}
 +9\bigl(9^d-6^d-2d6^{d-1}\bigr),
\]
que se simplifica a \eqref{mg03:ley-dimensional}. La división por
\(9^d\), seguida de la resta de la media aditiva \(5\), da
\eqref{mg03:limite-dimensional}; el término
\((d+3)(2/3)^{d-1}\) tiende a cero.
\end{proof}

En dimensión dos reaparecen \(S_{+,2}=405\),
\(S_{\times,2}=459\) y el defecto total \(54\). El límite
\(4\) pertenece al defecto medio bajo aumento de dimensión. El
defecto total, el cociente conservado y la media dimensional son
lectores especificados de la aritmética APP; TRIT y TPK transportan
posteriormente esos datos con su orientación, módulo y residencia.
